\section{问题四：2026 年 9 月 11--17 日的投放策略与期望范围}
\label{sec:q4}
本节先为每个单元建立日效率、CPC、CTR 与展位的日历预测模型并用滚动回测检验其点预测与区间，再在问题三的分配模型上叠加不确定性：证明期望注册最大化问题与用期望乘子的确定性问题同解，从而以两层分配给出逐日策略，并用蒙特卡罗情景给出竞价（CPC）、展现量、展现位、点击量、浏览量与注册量的期望范围。

\subsection{预测对象、日历模型与滚动回测}
\label{sec:q4model}
对单元 $u$ 有投放的每一天定义六个目标：效率乘子 $\log\varepsilon_{u,t}=\log C_{u,t}-\gamma_u\log(S_{u,t}/\bar s_u)-\log\bar c_u$（$\bar s_u,\bar c_u$ 为活跃日均值，即式~\eqref{eq:elast} 的残差结构）、$\log\mathrm{CPC}_{u,t}$、$\operatorname{logit}\mathrm{CTR}_{u,t}$、$\operatorname{logit}p^{\mathrm{top}}_{u,t}$、$\operatorname{logit}p^{1}_{u,t}$ 与 $\log C_{u,t}$。每个目标 $y_{u,t}$ 用日历回归
\begin{equation}\label{eq:fcmodel}
y_{u,t}=\beta_{u,0}+\sum_{w\ne\text{周三}}\beta_{u,w}D^{w}_t+\sum_{m}\mu_{u,m}M^{m}_t+\delta_{u,H}H_t+\delta_{u,A}A_t+e_{u,t},\qquad e_{u,t}\sim(0,\sigma_u^2),
\end{equation}
（月份效应不设基准月，由岭正则识别）以岭正则（$10^{-3}$，不含截距）稳定估计。预测点在样本末 8.5 个月之后，近期水平不可用，因此 2026 年的预测标准差取 $\sqrt{\sigma_u^2+(\sigma^{m}_u)^2}$，其中 $\sigma^{m}_u$ 为 2025 年月份效应的横截面标准差。必须强调：$\sigma^{m}_u$ 度量的是模型已经拟合掉的年内季节性离散度，而这里把它\emph{借用}为年际漂移的标准差——后者在只有一年数据时不可识别，这是一条不可检验的建模约定（假设~\ref{as:stable}），它使部署区间比仅含残差的区间宽 \valQfourWidthInflationMedian{} 倍（中位数；最小 \valQfourWidthInflationMin、最大 \valQfourWidthInflationMax）。9 月效应取 2025 年 9 月的估计（该月无投放的单元取基准月）。点预测经 $\exp$ 或 $\operatorname{expit}$ 反变换，效率乘子的期望取对数正态均值 $\hat\varepsilon_{u,t}=\exp(\mu_{u,t}+\tfrac12\sigma^2)$。

模型用 2025 年第四季度的滚动起点回测检验\cite{hyndman2021}：起点为 10 月 9 日起每周一个（共 12 个），预测步长 7 天，训练集至少 40 天；对照两条基线——季节朴素（上周同一周几的值，区间取 7 阶差分的 10\%/90\% 分位）与 28 日均值（区间取近 28 日偏差的分位）；指标为变换尺度上的 MAE、RMSE、分位数 0.1 与 0.9 的 pinball 损失\cite{koenker1978,gneiting2007}以及 80\% 区间覆盖率。为使回测检验的区间与实际交付的区间是\emph{同一个对象}，回测同时评估两种日历回归区间：仅含残差标准差 $\sigma_u$ 的“残差区间”，以及实际部署的 $\sqrt{\sigma_u^2+(\sigma^{m}_u)^2}$ 宽度的“部署区间”。附录表~\ref{tab:backtest} 与图~\ref{fig:backtest} 给出结果：残差区间的 80\% 覆盖率为 \valBacktestCalendarCoverageMin\%--\valBacktestCalendarCoverageMax\%（接近名义水平），部署区间为 \valBacktestWideCoverageMin\%--\valBacktestWideCoverageMax\%（明显\emph{过覆盖}，即偏保守）；28 日均值的区间过窄（\valBacktestMeanCoverageMin\%--\valBacktestMeanCoverageMax\%），季节朴素的区间约 90\%。因此本文报告的 10\%--90\% 范围应理解为\textbf{保守区间}：在 7 天步长的回测上它覆盖了约九成而非八成的实现值；对 8.5 个月之后的目标点，这一保守性是有意保留的，但读者应知道多出的宽度来自不可检验的 $\sigma^{m}$ 约定而非回测校准。点预测方面，6 个目标中 5 个的 MAE 由 28 日均值取得（近期水平的持续性强于全年日历结构）；点击量上日历回归与季节朴素基本持平（MAE \valBacktestCalendarMaeClicks{} 对 \valBacktestNaiveMaeClicks），日历回归的 RMSE 与 pinball 略优。由于 2026 年 9 月的目标点没有可用的“近期水平”，本文采用日历回归给出点预测；这一取舍与代价在第~\ref{sec:evaluation} 节中说明。

\paperfigure[0.86\linewidth]{fig_backtest.pdf}{滚动回测：日历回归、季节朴素与 28 日均值对各目标的 MAE（左）与 80\% 区间覆盖率（右，虚线为名义水平 80\%）。}{fig:backtest}

\subsection{预算、不确定性模型与随机分配}
\label{sec:q4alloc}
预算按“不超过 2025 年的投入”取 2025 年同期（9 月 11--17 日）各单元的实际消费之和 $B_u$（合计 \valQfourBudget{} 元），同期无消费的 \valQfourSkippedUnits{} 个单元在主口径下预算为零、不安排投放（在 \texttt{result4.xlsx} 中以零行显式给出），其余 \valQfourUnits{} 个单元参与；按全年日均折算（$\sum_tS_{u,t}/365\times7$，合计 \valQfourBudgetAnnualAlt{} 元）作为备选口径并同样给出完整方案（MDR-0008、MDR-0012）。不确定性以四类冲击进入响应（假设~\ref{as:day}）：单元--日共同冲击 $\varepsilon_{u,t}=\exp(\mu_{u,t}+\sigma_{u}^{\mathrm{tot}}w_t)$，$w_t\sim N(0,1)$，$\sigma_{u}^{\mathrm{tot}}=\sqrt{\sigma_u^2+(\sigma^{m}_u)^2}$；关键词特异冲击 $\xi_i=\exp(\sigma^{\mathrm{kw}}_u\zeta_i-\tfrac12(\sigma^{\mathrm{kw}}_u)^2)$（均值 1，$\zeta_i\sim N(0,1)$ 与 $w_t$ 独立，$\sigma^{\mathrm{kw}}_u$ 取该单元 $\log\mathrm{CPC}$ 的残差标准差）；CTR 与上方位占比的 logit 日冲击；注册的泊松噪声。于是
\begin{equation}\label{eq:stoch}
\begin{gathered}
\text{点击}_i=q_i(x_i)\,\varepsilon_{u,t}\,\xi_i,\qquad
\mathrm{CPC}_i=\frac{x_i}{\text{点击}_i},\qquad
\text{展现}_i=\frac{\text{点击}_i}{\mathrm{CTR}_{u,t}},\\
\text{浏览}_i=d_i\,\text{点击}_i,\qquad
\text{注册}_i\sim\mathrm{Poisson}(\rho_i\,\text{点击}_i),
\end{gathered}
\end{equation}
展位由式~\eqref{eq:position} 在情景 CPC 下计算。关键词的“竞价”在附件中不可观测，本文以其平均点击成本 $\mathrm{CPC}_i$ 作为竞价水平的度量。

\begin{proposition}[确定性等价]\label{prop:ce}
若冲击 $\varepsilon_{u,t}$ 与 $\{\xi_i\}$ 相互独立、二者均与投放额 $x$ 无关且 $E[\xi_i]=1$，则期望注册最大化问题 $\max_x E\bigl[\sum_i\rho_i q_i(x_i)\varepsilon_{u,t}\xi_i\bigr]$ 与用期望乘子 $\hat\varepsilon_{u,t}=E[\varepsilon_{u,t}]$ 的确定性问题 $(\mathrm P_{u,t})$ 的最优解集重合。
\end{proposition}
\begin{proof}
期望与有限和交换；由独立性 $E[\varepsilon_{u,t}\xi_i]=E[\varepsilon_{u,t}]E[\xi_i]=\hat\varepsilon_{u,t}$，故对任意可行 $x$ 两个目标值相等，最优解集因而重合。
\end{proof}
式~\eqref{eq:stoch} 的实现按独立抽样构造（$w_t$ 与 $\zeta_i$ 分别独立抽取），故命题的条件被满足。
因此最优策略无需情景优化，蒙特卡罗只用于给出各量的期望范围；样本平均近似的期望值与确定性目标一致，是模拟正确性的内部检验。策略分两层求解：(i) 跨日分配 $\max\{\sum_{t}\hat\varepsilon_{u,t}X_t^{\gamma_u}:\ \sum_tX_t\le B_u,\ X_t\in\{0\}\cup[m,\bar X_u]\}$，$\bar X_u=\max\{1.5\max_{t\in\text{同期}}S_{u,t},\ B_u/7\}$，用算法~\ref{alg:twophase}（把“日”视为关键词，$s=c=1$、系数 $\hat\varepsilon_{u,t}$）求解；(ii) 日内分配即 $(\mathrm P_{u,t})$，预算取 $X_t$。两层共 \valQfourAuditGroups{} 个问题全部通过独立校验器。需要明确：内层值函数 $V_t(X_t)$ 只在日内上限与最小投放额都不起作用时才等于 $K_uX_t^{\gamma_u}$，本数据上上限经常起作用，因此两层分解\emph{不保证}联合最优。为给出可核查的界，本文把每个单元的“7 天 $\times$ 候选词”联合问题写成一个可分离凹规划（去掉最小投放额、上限取 $\max\{\bar x_i(B_u),m\}$ 这一最宽的盒子），用同一注水机器求解：它的可行集包含任何两层解，故其最优值是联合最优的上界。实测联合间隙为 \valQfourJointGapPct\%（\valQfourJointUnits{} 个单元合计），单个单元最大 \valQfourJointMaxGapPct\%，出现在预期注册仅 \valQfourJointMaxGapUnitRegs{} 人的最小单元 \valQfourJointMaxGapUnit；在预期注册不低于 10 人的 \valQfourJointBigUnits{} 个单元上最大为 \valQfourJointMaxGapBigPct\%。随后对每个单元--日按式~\eqref{eq:stoch} 生成 \valQfourScenarios{} 个情景，报告每个关键词、每个单元--日与合计的期望值与 10\%--90\% 分位。

\subsection{结果}
\label{sec:q4res}
表~\ref{tab:q4-daily} 给出逐日策略的汇总，表~\ref{tab:q4-units} 给出各单元的 7 日结果，图~\ref{fig:q4-plan} 与图~\ref{fig:q4-units} 给出期望与范围。7 天预算 \valQfourBudget{} 元（最优解用尽至 0.3 元以内），入选 \valQfourSelectedKeywords{} 个关键词--单元组合（\valQfourRows{} 条关键词--日记录；不同关键词 ID 数略少，差额来自第~\ref{sec:q1kpi} 节的跨单元重复投放），预期点击 \valQfourClicksMean{} 次（10\%--90\% 分位 \valQfourClicksLow--\valQfourClicksHigh）、注册 \valQfourRegsMean{} 人（\valQfourRegsLow--\valQfourRegsHigh）、展现 \valQfourImpMean{} 次（\valQfourImpLow--\valQfourImpHigh）、浏览 \valQfourViewsMean{} 次（\valQfourViewsLow--\valQfourViewsHigh）、平均 CPC \valQfourCpcMean{} 元（\valQfourCpcLow--\valQfourCpcHigh）、点击加权的预期展位指数 \valQfourPositionMean（\valQfourPositionLow--\valQfourPositionHigh）。$\pi$ 的支撑集只有 $\{1,2.5,4\}$，\valQfourPositionMean{} 是一个三点混合的加权平均而非名次区间：其背后的点击加权展现结构是上方位占比 \valQfourPtopPct\%（其中首位 \valQfourPfirstPct\%、上方位其余 \valQfourPmidPct\%）、上方位之外 $100-\valQfourPtopPct$\%。又因为 $\pi$ 把所有非上方位展现一律记为 4（实际名次可能远在 4 之后），它是真实期望名次的\emph{乐观下界}。题目所要求的六个量的期望与范围汇总于表~\ref{tab:q4-ranges}，其中“竞价”以平均点击成本 CPC 度量、“展现位”以式~\eqref{eq:position} 的排位指数 $\pi$ 度量（附件不含出价与真实排名，二者只能用可观测量代理，MDR-0011）。合计量的区间由各推广单元的情景分位数相加得到（单元内 7 天按同一情景对齐求和，单元之间按分位数相加；表~\ref{tab:q4-daily} 的逐日行同理为当日各单元之和）：各单元的情景独立抽取，跨单元对齐没有概率含义，分位数相加等价于假设单元间冲击完全正相关，给出区间宽度的保守（偏宽）一侧；CPC 与展位不可加，改用点击量加权平均各单元--日的分位数。蒙特卡罗期望注册 \valQfourRegsMean{} 人与确定性目标 \valQfourDetRegs{} 人相差 \valQfourRegsMcZ{} 个蒙特卡罗标准误（SE $\approx\valQfourRegsMcSe$ 人），与命题~\ref{prop:ce} 相容。跨日分配给 9 月 13、14 日（周日、周一）的预算明显低于工作日（表~\ref{tab:q4-daily}），与问题一中周日、周一效率最低的发现一致。表~\ref{tab:q4-daily} 的“单元数”在 9 与 10 之间波动，原因是跨日分配允许 $X_t=0$：效率最低的日子上个别单元不投放。表~\ref{tab:q4-ranges} 的 CPC 行是\emph{点击加权}口径（7 日总投放 / 7 日期望点击），不是情景平均的 CPC 期望；按后者（各单元--日情景 CPC 均值的点击加权）为 \valQfourCpcScenarioMean{} 元，两者的差由 Jensen 不等式解释。展现量的范围远宽于点击量，因为它同时承受 CTR 与效率乘子的冲击；CPC 的范围（约 $\pm$30\%）反映了竞价的逐日波动。

\papertable{tab_q4_ranges}{2026 年 9 月 11--17 日六个量的期望与 10\%--90\% 范围：左半为 7 日合计（点击量、展现量、浏览量、注册量）或点击加权平均（CPC、展位指数 $\pi$），相对宽度 = 区间宽度 / 期望；右半为单个关键词--日的中位数（按 \valQfourRows{} 个关键词--日等权，因而由长尾词主导，与点击加权的合计值不可直接比较）}{tab:q4-ranges}

\papertable{tab_q4_daily}{2026 年 9 月 11--17 日逐日策略汇总：投入、CPC、展现量、展位、点击量、浏览量、注册量的期望与 10\%--90\% 范围}{tab:q4-daily}

\papertable{tab_q4_units}{各推广单元的 7 日预算、弹性、候选词数、CPC 与点击量、注册量的期望和 10\%--90\% 范围}{tab:q4-units}

\paperfigure[0.9\linewidth]{fig_q4_plan.pdf}{2026 年 9 月 11--17 日逐日预期点击量、注册量与 CPC（点为期望，带为 10\%--90\% 分位）。}{fig:q4-plan}

\paperfigure[0.62\linewidth]{fig_q4_units.pdf}{各推广单元 7 天预期注册量及 10\%--90\% 范围（对数横轴）。}{fig:q4-units}

同期预算口径把 \valQfourBiggestUnitBudgetSharePct\% 的预算留给单元 \valQfourBiggestUnitId{}（2025 年 9 月中旬其他单元几乎停投），而注册率最高、CPC 最低的长尾单元 9811363528 只分到二百余元却贡献了一百多人的预期注册（表~\ref{tab:q4-units}）。若改按全年日均口径安排预算（\valQfourBudgetAnnualAlt{} 元，较同期口径高 \valQfourBudgetAltOverMainPct\%，并恢复 \valQfourSkippedUnits{} 个被停投的单元），确定性预期注册由 \valQfourRegsExpectedSamePeriod{} 人升至 \valQfourRegsExpectedAnnualAlt{} 人（$+\valQfourRegsGainAnnualAltPct\%$）。必须说明这一增幅同时包含三种变化：预算增加 \valQfourBudgetAltOverMainPct\%、参与单元由 \valQfourUnits{} 个增至 \valQfourAltUnits{} 个、以及跨日节奏的改变；因此它不是纯粹的“同额预算下的时间再分配收益”。即便如此，方向是清楚的：2025 年 9 月中旬效率最高的长尾单元被停投，这是投放节奏本身的问题。本文建议公司在 2026 年 9 月不要沿用同期的单元停投安排，并按该口径给出\emph{完整的逐日逐词方案}（\valQfourAltRows{} 行、\valQfourAltUnits{} 个单元、\valQfourAltKeywords{} 个关键词，投放合计 \valQfourAltSpend{} 元、预期点击 \valQfourAltClicks{} 次、预期注册 \valQfourAltRegs{} 人），见支撑材料 \texttt{forecast/plan\_annual\_alt.csv}。与 2025 年同期实际点击 \valQfourRefClicks{} 次相比，同样的钱在最优策略下预期点击 \valQfourClicksMean{} 次（$+\valQfourClickGainVsRefPct\%$），其中包含了日历模型对 9 月效率的估计，应结合区间理解。详细结果（期望值）按模板写入 \texttt{result4.xlsx}：\valQfourRows{} 条有投放的关键词--日记录，另加 \valQfourSkippedUnits{} 个同期口径下无预算的单元各 7 天的零行（投入金额与各期望量均为 0，表示当日不投放），使全部 \valUnitsCount{} 个推广单元都在结果文件中出现。每个关键词--日的 10\%、50\%、90\% 分位与上方位、首位占比在支撑材料 \texttt{forecast/plan.csv} 中给出；备选口径的完整方案见 \texttt{forecast/plan\_annual\_alt.csv}。
